\section{Conclusion}\label{sec:conclusion}

The \bxthree{} Framework advances a single, architecturally precise claim: that
autonomous systems must propagate every unresolved exception state to a named
human principal, bypassing all intermediate machine actors, as a compulsory
structural requirement rather than a runtime choice.  This mandatory bailout
property is not a feature that can be retrofitted to an existing agentic
architecture through better error handling, improved escalation heuristics, or
more comprehensive monitoring dashboards.  It is a structural property that
follows, as a matter of architectural necessity, from the functional separation
between Purpose, Bounds, and Fact layers --- and the Bailout Protocol is the
mechanism by which that consequence is made operative at runtime.

The framework's five named architectural pillars are co-equal and mutually
reinforcing.  Loop Isolation ensures that the spawn graph is a rooted directed
acyclic graph and that accountability attribution is never ambiguous.
Recursive Spawning fixes the accountability chain at the architectural level
and renders it immutable.  Spatial Firewall contains sub-agent failures within
execution domains and prevents lateral state corruption.  Sandbox Gate provides
the sole permitted cross-domain communication channel with mandatory authorization
verification.  The Bailout Protocol compels the final escalation for any
condition not resolved by the preceding four pillars, including conditions not
anticipated at design time --- the protocol's most critical property is that it
operates on exceptions the designers did not foresee, because those are the
exceptions that matter most.

The theoretical contribution of this paper is that the upstream accountability
guarantee --- that \bxthree{} systems \emph{fail upward into human consciousness,
never downward into autonomous chaos} --- is not a design choice but a logical
consequence of the three-layer architecture.  The practical contribution is
equally clear: that autonomous systems operating without such a mechanism are
structurally exposed to failure modes --- invisible exception absorption,
contingent escalation, and irreversible drift --- that cannot be eliminated by
incremental improvement to existing escalation heuristics.  They require a new
architectural primitive.  The \bxthree{} Framework is that primitive, and the
Bailout Protocol is the mechanism by which it is enforced.

The limitations analysis identifies the boundaries of this primitive as currently
specified.  Human response latency introduces a structural floor on exception
resolution time that constrains the framework's applicability to high-velocity
operational domains.  Adversarial substrate conditions introduce attack surfaces
that the framework cannot defend without a substrate integrity layer that falls
outside its current scope.  Scalability is bounded by the 1:1 anchor model, which
requires the anchor pool extension before the framework can operate in enterprise
multi-agent deployments.  Strategic misclassification by the Bounds Engine
creates an exploitation pathway that the framework's trigger-condition architecture
does not yet close.  Each of these limitations is addressed by a specific directed
research agenda --- formal verification, adaptive timeout mechanisms, distributed
anchor pools, and exploitation-resistant trigger conditions --- that constitutes
the most immediate extensions of this work.

The core claim of the \bxthree{} Framework --- that the upstream accountability
guarantee requires mandatory human escalation for all unresolved exception states,
enforced architecturally rather than by policy --- is not contingent on the
resolution of these open questions.  It is established by the structural
necessity of the three-layer architecture, and it holds whether or not the
directed research agenda is pursued.  What the open questions determine is
whether the framework's guarantees can be extended to high-stakes, high-velocity,
large-scale, and adversarial deployment contexts --- and the research program
delineated here maps a principled path from the current specification toward
that goal.
